\section*{Chapitre \ref{ch:s2:quantitative-investment-strategies} --- Stratégies d'investissement quantitatives}

\begin{solution}{exo:s2:quantitative-investment-strategies:1}
$0.007 / 0.10 = 0.07$ de ratio de Sharpe.
\end{solution}

\begin{solution}{exo:s2:quantitative-investment-strategies:2}
$1/\sqrt{10} \approx 0.32$~: un ratio de Sharpe de backtest de 0.5 n'est qu'à environ 1.6 erreurs types de zéro avant toute sélection.
\end{solution}

\begin{solution}{exo:s2:quantitative-investment-strategies:3}
$1 - 0.16/0.51 = 68.6~\%$ sur les moyennes~; la médiane de l'érosion après backtest de chaque indice est de 72.2~\%, car ce rapport est déformé par les indices dont les backtests sont courts.
\end{solution}

\begin{solution}{exo:s2:quantitative-investment-strategies:4}
Les rendements en production n'ont pas servi à choisir l'indice, si bien que leur erreur moyenne est nulle~; le backtest a été retenu parce qu'il était le plus élevé des vingt, et il
porte donc les erreurs positives les plus fortes. La sélection biaise la statistique sur laquelle elle opère, pas l'avenir.
\end{solution}

\begin{solution}{exo:s2:quantitative-investment-strategies:5}
Chaque règle, filtre et paramètre ajouté constitue un essai de plus, et le backtest de la stratégie complexe est le meilleur parmi davantage d'essais~; ses erreurs sont
plus grandes, et une part plus importante de son backtest n'est que bruit. La complexité peut aussi épouser des particularités d'une seule période qui ne se reproduisent pas.
\end{solution}

\begin{solution}{exo:s2:quantitative-investment-strategies:6}
Quiconque sait quand l'indice se rééquilibre et, à peu près, ce qu'il va négocier peut acheter avant lui et vendre avant lui~;
l'indice s'exécute alors à de moins bons prix que ne le suppose un backtest sur les cours de clôture historiques.
\end{solution}

\begin{solution}{exo:s2:quantitative-investment-strategies:7}
Le meilleur backtest passe de 0.26 avec 2 candidats à 0.76 avec 100~; le résultat en production après coûts passe de 0.07 à 0.24~; l'écart passe de 0.18 à 0.52. La
médiane d'érosion reste comprise entre 68~\% et 76~\%, car davantage de candidats permettent aussi de trouver de meilleures stratégies~: le résultat en production augmente lui aussi, mais beaucoup moins que le
backtest. L'écart de ratio de Sharpe, et non le pourcentage, est ce que la complexité gonfle ici.
\end{solution}

\begin{solution}{exo:s2:quantitative-investment-strategies:8}
Réussir le test indique que le meilleur backtest est peu susceptible de relever du pur hasard, et non qu'il est exact~: ceux dont le backtest atteignait 1.13 et la production 0.32 avant
coûts. Le biais de sélection demeure dans le backtest, et les frais ainsi que les fuites de données s'y ajoutent.
\end{solution}

\begin{solution}{pb:s2:quantitative-investment-strategies:1}
\begin{enumerate}
\item Une stratégie fondée sur des règles, publiée par une banque sous forme d'indice et accessible au moyen de swaps, de notes structurées ou de fonds~; l'\omterm{def:s2:quantitative-investment-strategies:wrapper}{enveloppe de stratégie} désigne la forme sous laquelle un client
la détient.
\item Les frais, la marge sur les swaps et les notes, et la compensation des couvertures de l'indice avec ses autres flux.
\item 2\,000 équipes, comptant chacune 20 candidats dont les ratios de Sharpe réels sont voisins de 0.10 (écart-type 0.15) et des erreurs corrélées à 0.5, des backtests sur dix ans, les trois meilleurs
lancés et suivis en production pendant cinq ans à 10~\% de volatilité.
\item La distribution des ratios de Sharpe réels, si bien que l'érosion médiane est proche des 73~\% rapportés pour les stratégies construites par des banques.
\item 0.07 de ratio de Sharpe pour 0.7~\% par an à 10~\% de volatilité.
\item Personne ne les rémunérait avant que l'indice n'existe, et un backtest exécute ses ordres à des prix historiques que personne n'a pu devancer.
\item En gérant les positions de l'indice en miroir du swap.
\item D'autres opérateurs négocient avant le rééquilibrage publié~; l'indice achète plus haut et vend plus bas.
\item La chute du ratio de Sharpe du backtest au ratio de Sharpe en production après coûts, pour l'ensemble des indices lancés.
\item Les trois meilleurs~: backtest 0.51, réel 0.22, en production 0.23 avant coûts et 0.16 après~; érosion médiane 72.2~\%. Meilleur unique~: 0.62, 0.26, 0.27, 0.20, 69.4~\%.
\item 36.6~\% ont un ratio de Sharpe négatif en production après coûts.
\item Le test est franchi par 5.9~\% des meilleurs indices des équipes~; ceux dont le backtest atteignait 1.13 et la production 0.32, les autres 0.59 et 0.27 avant coûts.
\item Les indices lancés affichent en backtest un ratio de Sharpe de 0.51 et en production 0.16 après coûts~: un écart de 0.35 et une érosion médiane de 72~\%.
\item Une chute médiane de 73~\% du ratio de Sharpe, du backtest à la production, sur 215 stratégies bancaires, soit plus de 30 points de plus pour les plus complexes.
\item Davantage de candidats font passer le meilleur backtest de 0.26 à 0.76 et le résultat en production de 0.07 à 0.24~: l'écart passe de 0.18 à 0.52.
\item Les rendements des prédicteurs publiés étaient inférieurs de 26~\% hors échantillon et de 58~\% après publication.
\item Comme meilleur parmi un nombre inconnu d'essais, avant frais~: demander combien ont été testés, déflater, déduire les coûts et chercher des années en production.
\item L'indice volatilité-carry~: les expositions à la volatilité actions figuraient parmi les plus robustes en production.
\item Ce chapitre donne à voir le surapprentissage du livre~7, mesuré sur des produits effectivement vendus.
\item Un accès à une stratégie empaquetée, dont les coûts sont chiffrés et les risques couverts, vendu sur la foi d'un backtest.
\end{enumerate}
\end{solution}

\begin{solution}{iq:s2:quantitative-investment-strategies:1}
Une stratégie fondée sur des règles, publiée sous forme d'indice par une banque, avec son règlement de marché et un agent de calcul. La banque qui verse la performance de l'indice au titre d'un swap se couvre
en détenant les positions de l'indice et en les rééquilibrant selon les règles.
\end{solution}

\begin{solution}{iq:s2:quantitative-investment-strategies:2}
Combien de variantes ont été testées et comment les règles ont été choisies~; si les coûts, les frais et une exécution réaliste sont inclus~; quel est le
ratio de Sharpe déflaté~; comment il s'est comporté sur des sous-périodes et hors échantillon~; et si un historique en production existe.
\end{solution}

\begin{solution}{iq:s2:quantitative-investment-strategies:3}
Étaler le rééquilibrage sur plusieurs jours, randomiser son calendrier dans une fenêtre donnée, négocier sur des fixings moins prévisibles que la clôture, ou le compenser
en interne~; chacune doit figurer dans le règlement de marché avant d'entrer en vigueur.
\end{solution}

\begin{solution}{iq:s2:quantitative-investment-strategies:4}
Glissement de la couverture par rapport au niveau de l'indice, risque de gap et de saut les jours de rééquilibrage, risque de contrepartie sur les swaps, risque de modèle dans les indices complexes,
risque de réputation lié à de mauvaises performances en production, et conflits entre le promoteur de l'indice et son desk de trading.
\end{solution}

\begin{solution}{iq:s2:quantitative-investment-strategies:5}
Que le niveau de l'indice suive exactement le règlement de marché, à partir de prix de référence auditables~; que les rééquilibrages et les corrections soient publiés comme le prévoient les règles~;
que les niveaux historiques ne soient jamais retraités en silence~; et que le calcul soit indépendant du desk de trading.
\end{solution}

\begin{solution}{iq:s2:quantitative-investment-strategies:6}
Le meilleur de $N$ erreurs normales indépendantes a une espérance d'environ $s\sqrt{2\ln N}$ pour $N$ grand (l'approximation des valeurs extrêmes)~; l'estimation en production
du candidat retenu a pour espérance sa vraie valeur, la même pour tous. Le backtest dépasse donc l'espérance en production d'environ
$s\sqrt{2\ln N}$~: 0.77 pour $N = 20$ et $s = 1/\sqrt{10}$, avant la correction, dont la convergence est plus lente.
\end{solution}
